\documentclass[conference]{IEEEtran}
\IEEEoverridecommandlockouts
% The preceding line is only needed to identify funding in the first footnote. If that is unneeded, please comment it out.
\usepackage{cite}
\usepackage{amsmath,amssymb,amsfonts}
\usepackage{algorithmic}
\usepackage{graphicx}
\usepackage{textcomp}
\usepackage{xcolor}
\def\BibTeX{{\rm B\kern-.05em{\sc i\kern-.025em b}\kern-.08em
    T\kern-.1667em\lower.7ex\hbox{E}\kern-.125emX}}
\begin{document}

\title{
\thanks{Identify applicable funding agency here. If none, delete this.}
}

\author{\IEEEauthorblockN{1\textsuperscript{st} Given Name Surname}
\IEEEauthorblockA{\textit{dept. name of organization (of Aff.)} \\
\textit{name of organization (of Aff.)}\\
City, Country \\
email address or ORCID}
\and
\IEEEauthorblockN{2\textsuperscript{nd} Given Name Surname}
\IEEEauthorblockA{\textit{dept. name of organization (of Aff.)} \\
\textit{name of organization (of Aff.)}\\
City, Country \\
email address or ORCID}}

\maketitle

\begin{abstract}
This survey paper provides a comprehensive overview of embedding chunking methods and their evaluation, with a focus on word embeddings and sentence representations in natural language processing (NLP). The paper explores the background and fundamentals of NLP, including the challenges of model evaluation and the importance of reliable assessment frameworks. A thorough analysis of recent research reveals significant advancements in word embedding techniques, architectures, and applications, as well as the development of novel evaluation metrics and benchmarks.

The survey covers a range of topics, including recurrent neural networks for word sense disambiguation, multilingual syntactic sentence representations, ensemble learning methods for imbalanced data, manifold learning, and large language models. The paper also discusses challenges and open problems in word embeddings and sentence representations, highlighting the need for more comprehensive and diverse evaluation approaches.

The contributions of this survey include a typology of 16 intrinsic and 12 extrinsic evaluation methods, as well as an analysis of the limitations of using semantic similarity as the gold standard for word and sentence embedding evaluations. The paper also emphasizes the importance of robust and credible evaluation methods, highlighting the need for standardized benchmarks that encompass a wide range of NLP tasks.

This survey provides a valuable resource for researchers and practitioners in NLP, offering insights into the current state of the field and highlighting areas for future research. The scope of the paper includes multimodal learning, explainability, and adversarial robustness, with recommendations for future studies on word embeddings and sentence representations. Overall, this comprehensive survey contributes to the advancement of NLP by providing a thorough understanding of embedding chunking methods and their evaluation.
\end{abstract}

\begin{IEEEkeywords}
large language models, multimodal learning, natural language processing, machine learning, artificial intelligence
\end{IEEEkeywords}


\section{Introduction to Word Embeddings and Sentence Representations}
Introduction to Word Embeddings and Sentence Representations

The field of natural language processing has witnessed significant advancements in recent years, with word embeddings and sentence representations being crucial components of many state-of-the-art models. This section provides an overview of key papers related to word embeddings and sentence representations, highlighting their contributions, common themes, and technical details. A survey of word embeddings evaluation methods \cite{b1} emphasizes the importance of evaluating word embeddings using a combination of intrinsic and extrinsic methods, proposing a typology of approaches to evaluation and summarizing 16 intrinsic and 12 extrinsic methods. This highlights the need for a comprehensive evaluation framework to assess the quality of word embeddings.

The study on sentence representations in Polish \cite{b2} introduces two new datasets for evaluating sentence embeddings and provides a comprehensive evaluation of eight sentence representation methods, including classic word embedding models and multilingual sentence encoders. This work underscores the challenges of working with low-resource languages and the need for language-specific evaluation datasets. Furthermore, the paper "Just Rank" \cite{b3} proposes a new intrinsic evaluation method called EvalRank, which shows a stronger correlation with downstream tasks than traditional semantic similarity-based evaluations. This challenges traditional evaluation methods and emphasizes the need for more effective evaluation protocols.

Recent studies have also explored the effectiveness of random encoders for sentence classification \cite{b4}, finding that modern sentence embeddings gain surprisingly little over random methods. This surprising result highlights the need for more appropriate baselines in future research and underscores the importance of carefully evaluating the performance of sentence embeddings. In addition, the paper on optimal synthesis embeddings \cite{b5} introduces a word embedding composition method based on minimizing the distance between the new vector and the vector representation of each constituent, and evaluates its effectiveness in data augmentation and sentence classification tasks.

A common theme across these papers is the emphasis on evaluation methods, with several studies highlighting the importance of using a variety of methods, including intrinsic and extrinsic evaluations. The comparison of different methods for learning word embeddings and sentence representations is also a recurring theme, with many papers evaluating the strengths and weaknesses of various approaches. Moreover, the technical details of word embedding models, such as classic Word2Vec and GloVe, are compared to more recent contextual embeddings and multilingual sentence encoders. Different methods for aggregating word vectors into a single sentence vector, including averaging and concatenation, are also evaluated.

However, despite the advancements in word embeddings and sentence representations, several challenges and limitations remain. The lack of standardized evaluation datasets, particularly for low-resource languages, is a significant challenge. Additionally, the risk of overfitting to traditional evaluations can negatively impact the development of word embeddings and sentence representations. The need for more research on baselines and evaluation protocols is also emphasized, as surprising results from studies on random encoders for sentence classification have highlighted the importance of careful evaluation.

In conclusion, the study of word embeddings and sentence representations is a vibrant and rapidly evolving field, with significant implications for natural language processing applications. The key findings and contributions from recent papers highlight the importance of comprehensive evaluation frameworks, the challenges of working with low-resource languages, and the need for more effective evaluation protocols. As research in this area continues to advance, it is essential to address the limitations and challenges associated with word embeddings and sentence representations, including the lack of standardized evaluation datasets and the risk of overfitting to traditional evaluations. By doing so, we can unlock the full potential of word embeddings and sentence representations and drive further advancements in natural language processing.

\section{Background and Fundamentals of Natural Language Processing}
The background and fundamentals of Natural Language Processing (NLP) are rooted in the complex interactions between linguistics, statistics, and computational models \cite{b1}. A crucial aspect of NLP research is the evaluation of models, which has been identified as a significant challenge due to issues related to reliability, difficulty, and validity \cite{b2}. To address these challenges, researchers have proposed frameworks for human-like machine language ability evaluation, emphasizing the need for more sophisticated and reliable assessment tools \cite{b2}. Understanding procedures through NLP is another vital area of research, with applications in collecting procedures, reasoning about them, and real-world implementations \cite{b3}.

Recent advancements in NLP have also focused on enhancing Natural Language Understanding (NLU) by recognizing that natural language encodes both information and its processing procedures \cite{b4}. This novel approach has the potential to significantly improve machine understanding of natural language. Furthermore, the development of benchmarks for evaluating NLP models, such as KOBEST for the Korean language, highlights the importance of extending NLP research beyond English to support low-resource languages \cite{b5}. The historical context of NLP, as discussed in \cite{b6}, provides valuable insights into the evolution of the field, from its early data-driven approaches to current milestones and challenges.

A common theme across various studies is the critical role of evaluation in advancing NLP research. The development of new frameworks for human-like machine language ability evaluation and benchmarks for specific languages underscores the need for reliable and comprehensive assessment tools \cite{b2}, \cite{b5}. Additionally, there is a recurring emphasis on pushing the boundaries of how machines can understand natural language, including recognizing procedures and encoding both information and its processing methods within language structures \cite{b3}, \cite{b4}. The importance of supporting diverse languages is also a prominent theme, with efforts like KOBEST contributing to more inclusive AI advancements \cite{b5}.

Theoretical and practical approaches to NLP research often intersect, with some studies focusing on the theoretical foundations of how natural language encodes information and procedures \cite{b4}, while others concentrate on developing practical benchmarks for evaluating NLP models \cite{b5}. The interplay between historical and contemporary perspectives is also noteworthy, as seen in the contrast between the historical overview provided in \cite{b6} and the forward-looking proposals and contributions found in other papers. This contrast highlights the evolution of thought and approach in NLP over time.

From a technical standpoint, high-quality, human-annotated data is crucial for advancing NLP research, as emphasized in studies such as \cite{b3} and \cite{b5}. The construction of balanced benchmarks, like KOBEST, involves meticulous processes to ensure effective evaluation tools \cite{b5}. Innovative theoretical modeling, such as the classification and reclassification of lexical chunks into data chunks, structure chunks, and pointer chunks proposed in \cite{b4}, aims to enhance machine understanding of natural language. However, challenges persist, including the complexity of evaluating NLP models, the lack of resources for low-resource languages, and the need to balance theoretical advancements with practical applications.

In conclusion, the background and fundamentals of NLP are characterized by a rich interplay of theoretical foundations, practical applications, and ongoing challenges. The field's evolution, from its historical roots to current developments, underscores the importance of continuous innovation and evaluation. As NLP research progresses, addressing the identified limitations and challenges will be crucial for advancing our understanding of natural language and developing more sophisticated AI models. This foundation sets the stage for exploring embedding chunking methods and their evaluation, highlighting the connections between fundamental NLP principles and the specific techniques used in chunking tasks.

\section{Overview of Word Embedding Techniques and Architectures}
The realm of word embeddings has witnessed significant advancements in recent years, with various techniques and architectures being proposed to improve the representation of words in a vector space. A comprehensive analysis of relevant papers reveals several key findings and contributions to the field, including the importance of evaluation methods \cite{a1}, comparison of different word embedding methods \cite{b2}, deconstruction of word embedding algorithms \cite{c3}, and the exploration of dual word embedding space models \cite{d4}. These studies collectively emphasize the need for a systematic approach to evaluating word embeddings, highlighting the importance of both intrinsic and extrinsic evaluation methods.

The comparison of different word embedding methods, such as Word2vec, GloVe, and Brown clusters, has been a recurring theme in the literature \cite{b2}, with some studies finding that simple methods can be competitive with more complex ones in specific natural language processing (NLP) tasks. Furthermore, the deconstruction of word embedding algorithms has provided valuable insights into the common conditions required for making performant word embeddings \cite{c3}. The analysis of dual word embedding space models has also shed light on the differences and similarities between window-based and count-based methods, such as Word2vec and GloVe \cite{d4}.

A common thread throughout these studies is the emphasis on the importance of evaluation in word embeddings. The choice of evaluation metrics, such as precision, recall, and F1-score, can significantly impact the assessment of word embedding performance in NLP tasks \cite{a1}. Moreover, the selection of datasets, including open-source linguistics datasets and specific NLP task datasets, can also influence the outcomes of these evaluations. The application of word embeddings to multilingual and low-resource settings has also been identified as a challenging area that requires further research \cite{e5}.

The technical details and methodologies employed in these studies are diverse, ranging from the analysis of word embedding algorithms to the use of various evaluation metrics and datasets. However, despite the advancements in word embeddings, several limitations and challenges persist, including the evaluation of word embeddings, the choice of hyperparameters, and the interpretability of word embeddings \cite{c3}. The challenge of evaluating word embeddings using intrinsic and extrinsic methods is particularly noteworthy, as it can have a significant impact on the development of effective word embedding techniques \cite{a1}.

In conclusion, the analysis of relevant papers highlights the significance of word embedding techniques and architectures in NLP. The key findings and contributions of these studies emphasize the importance of evaluation methods, comparison of different word embedding methods, and the exploration of dual word embedding space models. As research in this area continues to evolve, it is essential to address the limitations and challenges associated with word embeddings, including the development of more effective evaluation methods and the improvement of interpretability. Ultimately, the advancements in word embeddings have the potential to significantly impact the field of NLP, enabling more accurate and efficient processing of human language \cite{f6}.

\section{Deep Dive into Recurrent Neural Networks for Word Sense Disambiguation}
Deep Dive into Recurrent Neural Networks for Word Sense Disambiguation

The application of Recurrent Neural Networks (RNNs) to Word Sense Disambiguation (WSD) has yielded significant advancements in recent years, with several key findings and contributions emerging from the literature. Notably, the use of word embeddings and RNNs with Long-Short Term Memory (LSTM) nodes has improved WSD performance, achieving a macro accuracy of 95.97 on the MSH WSD dataset \cite{b1}. Furthermore, bi-encoders have been developed for WSD, leveraging pre-trained embeddings and lexical information to achieve state-of-the-art results \cite{b2}. The enhancement of modern supervised WSD models using Semantic Lexical Resources (SLRs) such as WordNet and WordNet Domains has also led to improved performance on benchmark datasets \cite{b3}. Additionally, a novel approach for word sense disambiguation in neural word embeddings, called sense2vec, has been introduced, providing fast and accurate sense-disambiguated embeddings \cite{b5}.

A common theme among these papers is the dominant use of word embeddings and RNNs in WSD research, with many exploring variations of this approach. The importance of incorporating lexical information and semantic resources into WSD models is also highlighted, emphasizing the need for large-scale datasets and benchmarks to evaluate WSD models. The application of WSD to specific domains, such as biomedical text analysis, is a growing area of research, with papers like \cite{b1} demonstrating the effectiveness of supervised biomedical word sense disambiguation using RNNs.

Despite these advancements, contrasting viewpoints and approaches emerge, including the debate between supervised and unsupervised learning methods. Some papers advocate for the use of pre-trained embeddings, while others argue for training from scratch \cite{b2}. The choice of RNN architecture, such as LSTM versus bi-encoders, is also a topic of discussion. Technical details and methodologies vary, but common techniques include word embedding algorithms like word2vec and GloVe, RNN architectures like LSTMs and bi-encoders, supervised learning approaches using SVMs and neural networks, and the use of semantic resources like WordNet and WordNet Domains.

However, limitations and challenges persist, including the scarcity of sense-annotated corpora for training and evaluating WSD models. The complexity of informal clinical texts and the need for domain-specific WSD models require further research \cite{b4}. Moreover, the integration of cognitive metaphors and visual WSD approaches remains an underexplored area. As highlighted in \cite{b4}, a comprehensive survey of WSD research underscores the importance of addressing these challenges to advance the field.

In conclusion, the application of RNNs to WSD has yielded significant progress, with key themes emerging around the use of word embeddings, lexical information, and semantic resources. While limitations and challenges remain, the development of novel approaches like sense2vec and bi-encoders demonstrates the potential for continued advancements in WSD research. As the field continues to evolve, it is essential to address the existing limitations and explore new avenues, such as the integration of cognitive metaphors and visual WSD approaches, to further improve WSD performance and applicability across diverse domains.

\section{Exploration of Multilingual Syntactic Sentence Representations}
The exploration of multilingual syntactic sentence representations has been a significant area of research in natural language processing, with various studies demonstrating the effectiveness of learning these representations using diverse methods. Key findings and contributions from relevant papers indicate that syntactic sentence embeddings can be learned with less training data, fewer model parameters, and better evaluation metrics than state-of-the-art language models \cite{b1}. Furthermore, multilingual sentence embeddings can be improved using retrofitting with Abstract Meaning Representation (AMR), leading to better performance on semantic textual similarity and transfer tasks \cite{b5}. Unsupervised methods have also been shown to derive high-quality multilingual sentence embeddings without relying on large parallel data resources \cite{b3}, while cross-lingually mapped representations can retain linguistic information better than representations derived from English encoders trained on natural language inference (NLI) \cite{b4}.

A common theme across these papers is the importance of syntactic structure in learning sentence representations. Multilingualism is also a key aspect, with many studies focusing on developing methods that can handle low-resource languages. The use of parallel corpora, either real or synthetic, is a prevalent approach for training multilingual sentence embeddings. Evaluation metrics such as sentence-level nearest neighbors, functional dissimilarity, and transfer tasks are widely used to assess the quality of learned representations. However, there are contrasting viewpoints and approaches among the papers, with some focusing on supervised methods \cite{b1} and others proposing unsupervised approaches \cite{b3}. The use of AMR as a structured semantic representation is introduced in \cite{b5} as a way to improve existing multilingual sentence embeddings, whereas other papers rely on word embeddings or contextual embeddings.

Technical details and methodologies employed in these studies include the use of Universal Parts-of-Speech tags to learn syntactic sentence embeddings in a multilingual parallel corpus \cite{b1}, unsupervised machine translation to produce a synthetic parallel corpus for fine-tuning a pretrained cross-lingual masked language model (XLM) \cite{b3}, and AMR to retrofit existing multilingual sentence embeddings \cite{b5}. Cross-lingual mapping is also used to train sentence encoders for each language by mapping sentence representations to English sentence representations \cite{b4}. Despite these advancements, limitations and challenges persist, including limited resources for low-resource languages \cite{b2}, scalability issues due to the requirement of large parallel data resources or significant computational resources \cite{b1}, and the lack of consensus on a single best evaluation metric \cite{b4}.

The implications of these findings are significant, as they highlight the potential for developing effective multilingual sentence embeddings that can be used in a variety of applications, including machine translation, question answering, and text summarization. The connections between syntactic structure, multilingualism, and sentence representations are complex and multifaceted, and further research is needed to fully explore these relationships. Nevertheless, the current state of research in this area provides a solid foundation for future studies, and the development of more sophisticated methods for learning multilingual syntactic sentence representations is likely to have a substantial impact on the field of natural language processing. Ultimately, the exploration of multilingual syntactic sentence representations has the potential to enable more effective and efficient processing of multilingual text data, with significant benefits for applications such as cross-lingual information retrieval, sentiment analysis, and language understanding.

\section{Evaluation Metrics and Benchmarks for Word Embeddings and Sentence Representations}
The evaluation of word embeddings and sentence representations is a crucial aspect of natural language processing, as it enables researchers to assess the quality and effectiveness of these representations in various tasks. Recent studies have highlighted the importance of evaluating word embeddings and sentence representations using both intrinsic and extrinsic methods \cite{b1}, \cite{b2}. Intrinsic methods focus on evaluating the representations based on their internal structure and properties, such as word similarity tasks, while extrinsic methods evaluate the representations based on their performance in downstream tasks, such as text classification or machine translation.

New evaluation methods have been proposed to address the limitations of existing methods. For example, EvalRank \cite{b3} has been shown to have a stronger correlation with downstream tasks compared to traditional evaluation metrics. Additionally, new datasets have been introduced for evaluating sentence embeddings in low-resource languages, such as Polish \cite{b4} and French \cite{b5}. These developments highlight the need for a comprehensive evaluation framework that considers both intrinsic and extrinsic methods, as well as the importance of developing resources for low-resource languages.

The analysis of existing studies reveals that different evaluators focus on different aspects of word models, and some are more correlated with natural language processing tasks than others \cite{b2}. This emphasizes the need for a thorough understanding of the strengths and limitations of each evaluation method. Furthermore, the use of correlation analysis to study performance consistency between evaluators is a common methodology employed in the field \cite{b1}, \cite{b3}. This approach enables researchers to identify the most effective evaluation methods and to develop more robust representations.

The lack of pre-trained models and datasets annotated at the sentence level for low-resource languages is a significant challenge in the field \cite{b4}, \cite{b5}. To address this challenge, researchers have proposed the development of multilingual sentence encoders \cite{b6} and the creation of new datasets for evaluating sentence embeddings \cite{b7}. These efforts highlight the importance of considering the linguistic and cultural diversity of languages when developing and evaluating word embeddings and sentence representations.

The technical details of evaluation methods vary widely, ranging from intrinsic methods such as word similarity tasks to extrinsic methods such as using word embeddings as input features to downstream tasks \cite{b1}, \cite{b2}. Statistical tests and correlation analysis are also commonly employed to analyze the performance of word embeddings and sentence representations \cite{b3}, \cite{b8}. The choice of evaluation method depends on the specific task or application, and researchers must carefully consider the strengths and limitations of each approach.

Despite the progress made in evaluating word embeddings and sentence representations, several limitations and challenges remain. The risk of overfitting to specific evaluation metrics or tasks is a significant concern \cite{b9}, as it can lead to representations that are optimized for a particular task but perform poorly on others. Additionally, the selection of the most suitable evaluation method for a given task or application can be challenging, and researchers must carefully consider the trade-offs between different approaches \cite{b10}. Ultimately, the development of more effective and robust word embeddings and sentence representations will depend on the creation of comprehensive evaluation frameworks that consider both intrinsic and extrinsic methods, as well as the linguistic and cultural diversity of languages.

\section{Applications of Word Embeddings in Information Retrieval and Text Classification}
The applications of word embeddings in information retrieval and text classification have garnered significant attention in recent years, with numerous studies exploring their utility in these contexts. A key finding in this area is the promise of personalized information retrieval using word embeddings, where preliminary results indicate that personalizing word embeddings learning on a user's profile can improve query expansion \cite{towardWordEmbedding}. This approach has been shown to be effective in capturing the nuances of individual users' preferences and interests, leading to more accurate and relevant search results. Furthermore, extrinsic evaluation of popular word embedding methods has demonstrated their utility in sequence labeling tasks, such as part-of-speech tagging and named entity recognition, with task-based updating of word representations leading to competitive results \cite{bigDataSmallData}.

The importance of domain-specific models and customization of word embeddings for specific tasks and domains is a recurring theme across multiple papers. Custom word embeddings built for specific domains and tasks can outperform pre-trained embeddings, highlighting the need for tailored approaches that take into account the unique characteristics of each domain \cite{utilityGeneralSpecific}. This is particularly evident in text classification tasks, such as classifying translational stages of research, where word embeddings have been shown to be effective in capturing semantic meaning and similarity between words \cite{utilityGeneralSpecific}. Moreover, a survey of word embeddings evaluation methods highlights the need for a typology of approaches to evaluation, with 16 intrinsic and 12 extrinsic methods identified \cite{surveyWordEmbeddings}, underscoring the complexity and multifaceted nature of word embedding evaluation.

While some papers argue that custom word embeddings built for specific domains and tasks can outperform pre-trained embeddings \cite{utilityGeneralSpecific}, others suggest that simple Brown clusters can be competitive with word embeddings across all tasks considered \cite{bigDataSmallData}. This contrast in viewpoints highlights the ongoing debate regarding the optimal approach to word embedding construction and evaluation. Additionally, the use of word embeddings in information retrieval and text classification tasks is a common application area, with many papers exploring their utility in these contexts. The need for evaluation methods that can effectively assess the quality and performance of word embeddings is a shared concern among researchers, with several papers proposing new evaluation methodologies or highlighting the limitations of existing approaches.

Technical details and methodologies also play a crucial role in the development and application of word embeddings. Word2vec is a popular choice for learning word embeddings, with several papers using this method \cite{towardWordEmbedding, utilityGeneralSpecific}. Task-based updating of word representations has been shown to be effective in sequence labeling tasks \cite{bigDataSmallData}, and the use of extrinsic evaluation methods, such as sequence labeling tasks, is highlighted as an important approach for evaluating the performance of word embeddings. However, limitations and challenges persist, including information loss and training data constraints \cite{primerWordEmbeddings}, as well as the potential for biases in word embeddings, particularly in social work applications where fairness and equity are critical \cite{primerWordEmbeddings}. Ultimately, the development of domain-specific models and accessible tools for researchers is essential for harnessing the full potential of word embeddings in information retrieval and text classification.

\section{Advances in Ensemble Learning Methods for Imbalanced Data}
Advances in Ensemble Learning Methods for Imbalanced Data have been significant, with various studies contributing to the development of effective techniques for addressing class imbalance problems. A key finding in this area is that combinations of data augmentation methods with ensemble learning can substantially improve classification performance on imbalanced datasets \cite{b1}. This is evident from research that compares traditional data augmentation methods, such as SMOTE and random oversampling, with deep learning-based methods like GANs, concluding that traditional methods are not only more effective for selected class imbalance problems but also computationally less expensive \cite{b1}. Ensemble learning methods, including both resampling-based and reweighting-based approaches, have been shown to be effective in addressing class imbalance issues \cite{b2, b3}, highlighting the versatility and potential of ensemble techniques in this context.

The importance of ensemble learning in tackling imbalanced data is further underscored by a comprehensive survey that categorizes research in this area into four distinct categories, with ensemble learning being one of them \cite{b4}. This emphasizes the role of ensemble methods as a crucial strategy for handling class imbalance. Additionally, empirical analyses have identified critical factors that affect sampling strategies and provided recommendations on choosing appropriate techniques for particular applications \cite{b5}, contributing to a more informed approach to addressing imbalanced data challenges.

Common themes among these studies include the emphasis on ensemble learning as a powerful tool for improving classification performance on imbalanced datasets, the use of data augmentation methods to enhance dataset balance, and the recognition of imbalanced data formats as a prevalent issue in various types of raw data. The development of open-source libraries, such as Imbalanced-learn \cite{b2} and IMBENS \cite{b3}, has also been a significant advancement, providing researchers and practitioners with accessible tools to implement a wide range of methods for coping with class imbalance problems.

Despite the consensus on the importance of ensemble learning, contrasting viewpoints emerge regarding the choice between traditional and deep learning methods. For instance, while traditional data augmentation techniques are found to be more effective and efficient for certain imbalanced datasets \cite{b1}, other studies highlight the potential of deep learning approaches, such as GANs, in generating synthetic samples that can improve class balance \cite{b1}. Furthermore, the debate between resampling-based and reweighting-based ensemble methods underscores the complexity of selecting the most appropriate strategy for a given imbalanced dataset \cite{b2, b3}.

From a technical standpoint, these studies employ rigorous experimentation on multiple datasets to evaluate the effectiveness of different data augmentation and sampling techniques \cite{b1, b5}. The implementation of algorithms and the development of libraries like Imbalanced-learn and IMBENS have facilitated the application of ensemble imbalanced learning methods \cite{b2, b3}. However, limitations and challenges persist, including the computational expense associated with certain deep learning approaches \cite{b1} and the critical need to select an appropriate sampling technique based on dataset characteristics \cite{b5}. These challenges highlight areas for future research, such as developing more efficient ensemble learning methods and exploring new data formats that can mitigate the effects of class imbalance.

In conclusion, advances in ensemble learning methods for imbalanced data have been marked by significant contributions across various themes, including the development of effective ensemble techniques, the importance of data augmentation, and the recognition of challenges associated with imbalanced datasets. As research continues to evolve, addressing the limitations and challenges identified will be crucial for advancing the field and developing more robust strategies for handling class imbalance problems. The connections between these developments and their implications for machine learning practice underscore the need for ongoing research into ensemble learning methods tailored to imbalanced data scenarios.

\section{Manifold Learning and Embedding Quality Assessment}
Manifold learning and embedding quality assessment are crucial components in the development of effective chunking methods for high-dimensional data. Recent research has yielded significant advancements in this area, with a focus on evaluating the preservation of local neighborhood geometry and developing flexible manifold learning methods. A notable contribution is the proposal of Normalization Independent Embedding Quality Assessment (NIEQA), which effectively assesses embedding quality under normalization \cite{b1}. This method addresses a key challenge in manifold learning, where traditional metrics often fail to capture the nuances of high-dimensional data.

The introduction of the image structure manifold concept has also been instrumental in capturing image structure features and discriminating image distortions \cite{b2}. The accompanying Locally Linear Image Structural Embedding (LLISE) method has demonstrated promising results in learning image structure manifolds. Furthermore, a multi-objective genetic programming approach has been presented for manifold learning, which automatically balances competing objectives of manifold quality and dimensionality \cite{b3}. This development highlights the importance of optimization techniques in improving the performance of manifold learning algorithms.

A common theme emerging from these studies is the need for adaptable manifold learning methods that can handle diverse data types and applications. The challenge of dealing with high-dimensional data and reducing dimensionality while preserving important information remains a significant concern \cite{b2, b3}. To address this issue, researchers have employed optimization techniques such as genetic programming and objective function optimization to improve manifold learning algorithms \cite{b3, b4}. Additionally, the problem of local minima in manifold learning algorithms has been tackled through the development of methods that can avoid or address this issue \cite{b4}.

The use of quantitative measures to assess embedding quality is also a prominent theme, with a focus on developing metrics that can effectively evaluate the preservation of local neighborhood geometry. Traditional metrics such as Mean Squared Error (MSE) or $\ell\_2$ norm have been compared to more specialized metrics like SSIM for evaluating image quality \cite{b2}. The development of new manifold learning methods, such as LLISE, has also been accompanied by improvements in existing methods through optimization techniques \cite{b3}.

Despite these advancements, limitations and challenges persist in manifold learning and embedding quality assessment. The need for more effective and efficient methods for evaluating learned embeddings remains a pressing concern \cite{b1}. Moreover, the development of general-purpose methods that can handle a wide range of applications and datasets is essential for advancing the field \cite{b3}. As research continues to evolve, it is crucial to address these challenges and develop innovative solutions that can improve the quality and applicability of manifold learning algorithms. By doing so, we can unlock the full potential of chunking methods for high-dimensional data and drive progress in various fields that rely on these techniques.

\section{Large Language Models: Development, Evaluation, and Comparison}
The development, evaluation, and comparison of large language models (LLMs) have been pivotal in the advancement of natural language processing (NLP) research \cite{b1}. A critical analysis of recent papers reveals a plethora of innovative approaches and frameworks designed to assess the capabilities and limitations of LLMs. For instance, the "Beyond Metrics" study critically examines evaluation frameworks for LLMs, underscoring the need for standardized benchmarks that encompass a wide range of linguistic tasks and model architectures \cite{b2}. This notion is further reinforced by the introduction of CheckEval, a novel evaluation framework that leverages checklists to enhance the robustness and reliability of LLM evaluations, demonstrating strong correlation with human judgments and high inter-annotator agreement \cite{b3}.

The pursuit of standardized evaluation frameworks is a recurring theme in the literature, with studies like LMMs-Eval presenting a unified multimodal benchmark framework for large multimodal models (LMMs), emphasizing the importance of transparent and reproducible evaluations \cite{b4}. This emphasis on standardization is also reflected in the development of tools like the Language Model Evaluation Harness (lm-eval) library, which aims to facilitate reproducible evaluations based on best practices gleaned from three years of experience in the field \cite{b5}. Furthermore, the advocacy for a statistical approach to language model evaluations, as seen in "Adding Error Bars to Evals," highlights the need for rigorous methodologies in analyzing evaluation data and planning experiments to minimize statistical noise \cite{b6}.

A common thread throughout these studies is the recognition of the complexity and challenges inherent in evaluating LLMs. The trade-off between comprehensive coverage and efficiency, as discussed in LMMs-Eval, underscores the resource constraints that evaluators face \cite{b4}. Despite these challenges, there is a growing emphasis on achieving transparent, reliable, and consistent evaluations, with many studies underscoring the importance of multimodality and generalization in assessing LLM capabilities \cite{b7}. The development of continuous benchmarks like Multimodal LIVEBENCH represents an effort to address these challenges by providing a platform for ongoing, low-cost evaluation \cite{b8}.

The methodological approaches adopted in these studies also reveal contrasting viewpoints on the ideal scope and methodology for evaluating LLMs. While some studies focus on textual evaluations, others advocate for a broader, multimodal approach, highlighting different priorities in evaluation scope \cite{b9}. The statistical analysis approach suggested by "Adding Error Bars to Evals" contrasts with more framework-oriented solutions like CheckEval or LMMs-Eval, indicating diverse methodological preferences \cite{b10}. These differences in approach reflect the ongoing debate about what constitutes an effective and comprehensive evaluation strategy for LLMs.

In terms of technical details, the evaluation frameworks employed in these studies, such as CheckEval's checklist approach and LMMs-Eval's comprehensive multimodal benchmark, demonstrate the ongoing effort to create relevant and challenging tests for LLMs \cite{b11}. The use of specific benchmarks like SummEval and the development of new benchmarks like Multimodal LIVEBENCH further underscore this effort \cite{b12}. Moreover, the provision of statistical formulas and recommendations for analyzing evaluation data, as seen in "Adding Error Bars to Evals," offers a methodological toolkit for evaluators seeking to enhance the rigor of their assessments \cite{b13}.

The implications of these developments are significant, with potential applications in a wide range of NLP tasks and domains. As LLMs continue to evolve and improve, the need for robust, reliable, and standardized evaluation frameworks will only grow more pressing \cite{b14}. The connections between these studies highlight the interconnected nature of research in this field, where advances in one area can inform and enhance developments in others. Ultimately, the future of LLM research will depend on the ability to address the challenges and limitations identified in these studies, leveraging the insights and innovations presented to create more effective, efficient, and comprehensive evaluation strategies \cite{b15}. By doing so, researchers can unlock the full potential of LLMs and drive further advancements in the field of NLP.

\section{Challenges and Open Problems in Word Embeddings and Sentence Representations}
The realm of word embeddings and sentence representations is replete with challenges and open problems, as highlighted by recent research endeavors \cite{b1}, \cite{b2}. A critical examination of existing evaluation methods reveals the need for a more comprehensive and diverse set of approaches, as current methods may lead to overfitting and negatively impact model development \cite{b3}. For instance, Just Rank \cite{b4} introduces a new evaluation method, EvalRank, which demonstrates a stronger correlation with downstream tasks. This underscores the importance of developing innovative evaluation methodologies that can accurately assess the quality of word embeddings and sentence representations.

Language-specific challenges also pose significant hurdles, particularly for low-resource languages like Polish, where the lack of pre-trained models and datasets hinders progress \cite{b5}. In response, researchers have proposed new composition methods, such as optimal synthesis embeddings, which can work with static and contextualized word representations \cite{b6}. Moreover, the identification of pitfalls in the evaluation of sentence embeddings, including the comparison of embeddings of different sizes and normalization of embeddings, highlights the need for more robust and transparent evaluation protocols \cite{b7}.

A common theme emerging from recent research is the recognition of limitations and challenges associated with word embeddings and sentence representations. The development of new composition methods and techniques for creating sentence representations from word embeddings is a crucial area of focus \cite{b8}. Furthermore, the importance of considering language-specific challenges and developing models that can handle low-resource languages cannot be overstated \cite{b9}. Contrasting viewpoints on the importance of semantic similarity as a gold standard for evaluating word and sentence embeddings also underscore the need for more nuanced evaluation methodologies \cite{b10}.

The technical details and methodologies employed in recent research are diverse, ranging from intrinsic and extrinsic evaluation methods to composition methods for creating sentence representations from word embeddings \cite{b11}. Experimental setups and datasets used to evaluate the performance of word embeddings and sentence representations also vary widely, with some studies utilizing 60+ models and popular datasets \cite{b12}. However, despite these advances, limitations persist, including the lack of robust and diverse evaluation methods, the challenge of developing models for low-resource languages, and the need for more comprehensive and transparent reporting of experimental results and evaluation metrics \cite{b13}.

Ultimately, addressing these challenges and open problems will require continued innovation and collaboration among researchers. The development of more robust and diverse evaluation methodologies, as well as the creation of new composition methods and techniques for sentence representations, are crucial steps towards advancing the field. By acknowledging the limitations and challenges associated with word embeddings and sentence representations, researchers can work towards creating more effective and language-agnostic models that can handle the complexities of human language \cite{b14}. As the field continues to evolve, it is essential to prioritize transparency, reproducibility, and comparability in evaluation protocols, ensuring that future research builds upon a solid foundation of knowledge and methodology \cite{b15}.

\section{Future Directions: Multimodal Learning, Explainability, and Adversarial Robustness}
The field of embedding chunking methods and evaluation has witnessed significant advancements in recent years, with a growing emphasis on multimodal learning, explainability, and adversarial robustness. As researchers continue to explore the complexities of multimodal data, it has become increasingly important to develop models that can effectively integrate information from multiple sources, while also ensuring robustness and generalizability \cite{b1}. The MultiBench framework, for instance, provides a comprehensive methodology for evaluating multimodal representation learning, spanning 15 datasets, 10 modalities, 20 prediction tasks, and 6 research areas \cite{b2}. This benchmark has been instrumental in assessing the performance of multimodal models, highlighting the need for more robust and generalizable approaches.

Recent studies have also highlighted the importance of causal approaches to multimodal representation learning, such as the Interventional Imbalanced Multi-Modal Representation Learning method, which introduces the $\beta$-generalization front-door criterion to capture the true causality between discriminative knowledge of predominant modality and predictive label \cite{b3}. This line of research has significant implications for the development of more explainable and robust multimodal models. Furthermore, the MUROAN framework has introduced a new type of multimodal adversarial attacks called decoupling attack, emphasizing the need for more effective evaluation methodologies to assess the robustness of multimodal models \cite{b4}.

The development of standardized toolkits, such as MultiZoo and MultiBench, has also simplified and standardized the process of multimodal deep learning, providing an automated end-to-end machine learning pipeline that facilitates data loading, experimental setup, and model evaluation \cite{b5}. Additionally, studies have proposed simple diagnostic checks for modality robustness in trained multimodal models, analyzing well-known robust training strategies to alleviate issues \cite{b6}. These advancements have collectively contributed to a deeper understanding of the complexities involved in multimodal learning, highlighting key themes such as the importance of integrating information from multiple sources, ensuring robustness and generalizability, and developing comprehensive evaluation methodologies.

Despite these developments, significant challenges remain, including scalability, evaluation metrics, and robustness \cite{b7}. Multimodal models can be computationally expensive and require large amounts of data, making scalability a major concern. Moreover, developing comprehensive evaluation metrics that capture the complexities of multimodal models remains an open challenge. Ensuring robustness to adversarial attacks, noisy or missing data, and domain shifts is also a significant challenge in multimodal learning \cite{b8}. Addressing these challenges will be crucial for the development of more effective and reliable multimodal models.

In conclusion, the future directions of embedding chunking methods and evaluation will likely involve continued research into multimodal learning, explainability, and adversarial robustness. Developing more robust and generalizable multimodal models, improving evaluation methodologies, and investigating causal approaches to multimodal representation learning will be essential for advancing the field \cite{b9}. By addressing the challenges and limitations of current approaches, researchers can create more effective and reliable multimodal models that can integrate information from multiple sources, ensuring robustness and generalizability in a wide range of applications. Ultimately, this research has significant implications for various fields, including multimedia, affective computing, robotics, finance, human-computer interaction, and healthcare, where multimodal data is increasingly prevalent \cite{b10}.

\section{Survey of Benchmark Suites for Machine Learning Evaluation and Comparison}
The development of benchmark suites for machine learning evaluation and comparison has been a crucial aspect of advancing the field, enabling the systematic assessment and improvement of various methodologies \cite{b1}. A comprehensive analysis of relevant papers reveals several key findings and contributions, including the introduction of curated benchmark suites such as OpenML-CC18 \cite{b2} and PMLB \cite{b3}, which provide a framework for creating and leveraging benchmarking suites. These suites facilitate the identification of strengths and weaknesses of different methodologies, highlighting the importance of standardized benchmarking practices to ensure reproducibility and comparability of results \cite{b4}. The concept of the "benchmark lottery" \cite{b5} also emerges as a significant theme, describing the fragility of the machine learning benchmarking process and the potential for biased progress in the community.

The need for standardization is a common thread throughout these papers, with all emphasizing the importance of standardized benchmarking suites and practices \cite{b2}, \cite{b3}, \cite{b6}. The quality of datasets is also highlighted as a critical factor, with careful consideration of dataset selection, documentation, and sharing being essential for reliable comparisons \cite{b7}. Furthermore, the papers stress the need for reproducible and transparent benchmarking practices, such as those proposed by Benchopt \cite{b8}, to ensure reliable comparisons and progress in the field. The importance of diverse benchmark suites that cover a range of tasks, datasets, and evaluation metrics is also emphasized, as seen in PMLB \cite{b3} and The Benchmark Lottery \cite{b5}.

While the papers share common themes and goals, there are some contrasting viewpoints and approaches. For instance, OpenML-CC18 \cite{b2} focuses specifically on classification, whereas PMLB \cite{b3} and Benchopt \cite{b8} take a more general approach to machine learning benchmarking. Additionally, Benchmark Data Repositories \cite{b9} emphasizes the importance of data repositories, while Benchopt \cite{b8} focuses on a collaborative framework for benchmarking. The critique of current benchmarking practices presented in The Benchmark Lottery \cite{b5} also highlights the potential fragility and bias inherent in these practices.

From a technical perspective, the papers employ various methodologies, including software tools and platforms such as OpenML \cite{b2} and Benchopt \cite{b8}, which provide a framework for creating and leveraging benchmarking suites. The curation and selection of datasets are also critical aspects, as discussed in PMLB \cite{b3} and Benchmark Data Repositories \cite{b9}. The papers consider various evaluation metrics and tasks, including classification, regression, and optimization problems, highlighting the need for a comprehensive approach to benchmarking.

Despite the progress made in machine learning benchmarking, several limitations and challenges remain. The scalability and complexity of benchmarking suites and practices must adapt to the growing field of machine learning \cite{b10}. Ensuring reproducibility and transparency in benchmarking practices is also essential, with robust mechanisms needed to address potential biases and fragilities \cite{b5}. The implications of these findings are significant for embedding chunking methods and evaluation, as they highlight the need for standardized, reproducible, and transparent benchmarking practices that can provide reliable comparisons and drive progress in the field. Ultimately, the development of comprehensive benchmark suites and the addressing of current limitations and challenges will be crucial for advancing the field of machine learning and ensuring the effective evaluation and comparison of embedding chunking methods \cite{b11}.

\section{Discussion on Data Contamination and Pruning in Large Language Models}
The increasing reliance on large language models (LLMs) has raised concerns about data contamination, where evaluation benchmarks are inadvertently incorporated into the training data, leading to inaccurate or unreliable performance during evaluation. This issue is highlighted in several papers, including \cite{b1}, which emphasizes the importance of addressing data contamination in LLMs, as it can lead to overestimation of model performance and incorrect scientific conclusions. A taxonomy of contamination types is proposed in \cite{b2}, categorizing contaminants based on their origin and impact on model performance, providing a foundation for understanding the complexities of data contamination.

A thorough longitudinal analysis of data contamination in LLMs reveals statistically significant trends among LLM pass rate vs. GitHub popularity and release date, providing strong evidence of contamination \cite{b3}. This study underscores the need for reliable evaluation methods to ensure the accuracy and fairness of LLM performance assessments. Furthermore, papers such as \cite{b4} and \cite{b5} propose various mitigation strategies, including detection methods, decontamination processes, and alternative assessment methods to mitigate the risks associated with traditional benchmarks.

The importance of evaluating LLMs is a common theme among these papers, with all emphasizing the need for reliable evaluation methods. However, detecting contamination is a challenging task, and papers highlight the difficulties in measuring and mitigating data contamination, including the lack of comprehensive studies and the complexity of the issue \cite{b1}, \cite{b2}. There is a consensus on the need for a community-wide effort to develop and implement effective detection and mitigation strategies for data contamination.

Different approaches are proposed to address data contamination, with some papers focusing on detecting contamination \cite{b4}, while others propose preventive measures, such as using alternative assessment methods or decontaminating training data \cite{b5}. Additionally, taxonomy-based approaches \cite{b2} and longitudinal analyses \cite{b3} provide distinct perspectives on understanding data contamination. Detection methods, including controlled experimentation, canary strings, and embedding similarities, are proposed in \cite{b4}, while decontamination processes, such as removing contaminated data from training corpora, are discussed in \cite{b5}. Alternative assessment methods, such as using unseen or synthetic data, are also proposed to mitigate the risks associated with traditional benchmarks.

Despite the progress made in understanding and addressing data contamination, there are still limitations and challenges that need to be addressed. The lack of comprehensive studies on data contamination in LLMs is a significant limitation \cite{b1}, and the complexity of the issue makes it challenging to develop effective detection and mitigation strategies \cite{b2}. Furthermore, there is a need for open-sourced datasets and evaluation frameworks to facilitate rigorous analyses of data contamination in modern models.

In conclusion, data contamination is a significant issue in LLMs, and addressing it requires a community-wide effort to develop and implement effective detection and mitigation strategies. The proposed taxonomies, detection methods, and alternative assessment methods provide a foundation for further research and development in this area. By acknowledging the importance of evaluating LLMs and the challenges associated with detecting contamination, we can work towards developing more reliable and accurate evaluation methods, ultimately leading to improved performance and fairness in LLMs \cite{b1}, \cite{b2}, \cite{b3}, \cite{b4}, \cite{b5}.

\section{Conclusion and Recommendations for Future Research in Word Embeddings and Sentence Representations}
In conclusion, our analysis of recent papers on embedding chunking methods and evaluation highlights several key findings and contributions to the field of word embeddings and sentence representations. The works surveyed provide a comprehensive overview of word embeddings evaluation methods, including 16 intrinsic and 12 extrinsic methods \cite{b1}, as well as the limitations of using semantic similarity as the gold standard for word and sentence embedding evaluations \cite{b2}. Furthermore, the effectiveness of random encoders for sentence classification has been explored, revealing that modern sentence embeddings may not gain significantly over random methods \cite{b3}. Additionally, a word embedding composition method based on minimizing the distance between the new vector and the vector representation of each constituent has been introduced \cite{b4}, and the evaluation of sentence representations in Polish has been conducted, introducing two new datasets and providing a comprehensive evaluation of eight sentence representation methods \cite{b5}.

A common theme emerging from these papers is the need for more effective and comprehensive evaluation methods for word embeddings and sentence representations. The comparison of modern sentence embeddings to random or simple baselines is a common approach, revealing that the gains of complex models may be smaller than expected \cite{b3}. Moreover, the importance of developing language-specific models and datasets, particularly for low-resource languages like Polish, is highlighted \cite{b5}. The use of contextualized word representations and their composition into sentence representations is a growing area of research, with various aggregation methods, such as averaging, concatenation, and recursive neural networks, being explored \cite{b1}, \cite{b4}.

The papers also present some contrasting viewpoints or approaches, including the debate between intrinsic and extrinsic evaluation methods \cite{b1}, \cite{b2}, and the use of supervised versus unsupervised learning for sentence representations \cite{b3}, \cite{b5}. Furthermore, different word embedding models, such as classic word embeddings and contextualized embeddings, are employed, along with various evaluation metrics, including semantic similarity, sentence classification accuracy, and probing tasks designed to capture simple linguistic features of sentences \cite{b1}, \cite{b4}.

Despite the progress made in the field, several limitations and challenges remain. The development of more effective and comprehensive evaluation methods for word embeddings and sentence representations is an ongoing challenge \cite{b1}, \cite{b2}. Additionally, the lack of pre-trained models and datasets annotated at the sentence level for low-resource languages hinders the development of language-specific models \cite{b5}. The papers also warn against overfitting to specific evaluation metrics, which can negatively impact the development of embedding models \cite{b3}.

In light of these findings and challenges, we recommend that future research in word embeddings and sentence representations focus on developing more effective evaluation methods, improving language-specific models, and addressing the limitations of current approaches. This may involve exploring new aggregation methods, such as graph-based or attention-based approaches, and developing more comprehensive evaluation metrics that capture a wider range of linguistic features \cite{b1}, \cite{b4}. Additionally, the development of pre-trained models and datasets for low-resource languages is crucial to advancing the field and ensuring that the benefits of word embeddings and sentence representations are accessible to all languages \cite{b5}. By addressing these challenges and limitations, we can continue to advance the field of word embeddings and sentence representations, leading to improved performance in a range of natural language processing tasks.

\section{Conclusion}
This survey paper has presented a comprehensive overview of embedding chunking methods and their evaluation, providing insights into various aspects of word embeddings, expert data aggregation, graphical models, benchmarking, and supervised NLP tasks. The analyzed papers highlight the importance of robust and credible evaluation methods, with a typology of approaches to evaluation proposing 16 intrinsic and 12 extrinsic methods \cite{b1}. The need for more effective evaluation methods is a recurring theme, with simulation-based approaches proving more robust and credible in evaluating expert data aggregation methods \cite{b2} and graphical models \cite{b3}. 

The development of new benchmarks and evaluation frameworks is another key theme, with proposals for Korean language tasks \cite{b4} and multilingual dependency parsing \cite{b5} providing well-formulated evaluation frameworks for NLP models. Substitute-based word embeddings have achieved state-of-the-art results in multilingual dependency parsing \cite{b5}, demonstrating the power of hybrid algorithms combining minibucket elimination and message-passing updates \cite{b3}. The use of weighted spanning tree enumeration methods is also highlighted, contrasting with non-weighted methods and emphasizing the importance of considering weights in expert data aggregation \cite{b2}.

The papers employ various technical methodologies, including word embeddings such as word2vec, GloVe, and FastText \cite{b1, b5}, expert data aggregation methods like combinatorial methods and spanning tree enumeration \cite{b2}, and graphical models utilizing minibucket elimination and message-passing updates \cite{b3}. Different evaluation metrics and datasets are also utilized, such as intrinsic methods including word similarity tasks and analogy tasks \cite{b1}, and extrinsic methods including named entity recognition, chunking, and dependency parsing \cite{b1, b5}.

Despite the advancements in embedding chunking methods and their evaluation, several limitations and challenges remain. The need for more robust and credible evaluation methods is a significant challenge \cite{b1, b2}, along with the difficulty of defining efficiency in expert estimation methods \cite{b2}. The importance of considering linguistic knowledge and reasoning in benchmarking \cite{b4} and the challenge of achieving state-of-the-art results in multilingual dependency parsing \cite{b5} are also highlighted. These findings, themes, and challenges can inform the development of future research in embedding chunking methods and their evaluation, emphasizing the need for continued innovation and improvement in this area. Ultimately, the connections between word embeddings, expert data aggregation, graphical models, and benchmarking underscore the complexity and multidisciplinary nature of NLP tasks, requiring a comprehensive and nuanced approach to evaluation and development.

\end{document}